\chapter{就活文書の作成}
\label{chap9}

\begin{chapterbox}
\textbf{【この章で学ぶこと】}
\begin{itemize}
\item AIを活用してガクチカ・ESを効果的に作成するワークフローを習得する
\item STAR法による経験の構造化技法を身につける
\item 研究内容を非専門家に分かりやすく伝える技法を理解する
\end{itemize}
\end{chapterbox}
\vspace{5mm}

\section{大学院生の就職活動とAI}
\index{しゅうしょくかつどう@就職活動}

\subsection*{M1の就活スケジュール}

大学院修士課程の学生にとって、就職活動は研究と並行して進めなければならない重要なタスクである。

\begin{table}[H]
\centering
\begin{tabular}{ll}
\toprule
時期 & 活動内容 \\
\midrule
M1 4月〜6月 & 自己分析、業界研究の開始 \\
M1 6月〜8月 & サマーインターンシップ応募・参加 \\
M1 9月〜11月 & 秋冬インターンシップ、OB/OG訪問 \\
M1 12月〜2月 & 本選考に向けたES準備、早期選考開始 \\
M1 3月〜 & 本選考エントリー開始（広報解禁） \\
M2 4月〜6月 & 面接・内定（選考解禁は6月1日が建前） \\
\bottomrule
\end{tabular}
\end{table}

\subsection*{企業が大学院生に求めるもの}

企業が大学院生を採用する際に注目するポイントは、大きく分けて2つある。

\begin{enumerate}
\item \textbf{専門性} --- 専門知識そのものだけでなく、「深く学ぶ力」の証明。研究を通じて培った分析力、論理的思考力。技術的な課題に粘り強く取り組む姿勢。
\item \textbf{汎用スキル} --- 課題発見力、仮説構築力、コミュニケーション力、プロジェクト管理力、主体性。
\end{enumerate}

重要なのは、企業は「研究の専門知識をそのまま使いたい」のではなく、「研究を通じて培われた思考プロセスと行動特性」を評価しているということである。

\subsection*{AIを就活で活用する意義と注意点}

\begin{itembox}[l]{AI活用の意義と注意点}
\textbf{意義}
\begin{itemize}
\item 自己分析の壁打ち相手として、客観的な視点を得られる
\item 文章の構造化・推敲を効率的に行える
\item 想定質問の生成により、面接準備の抜け漏れを減らせる
\item 専門用語の平易化など、表現の調整に役立つ
\end{itemize}

\textbf{注意点}
\begin{itemize}
\item AIが生成した文章をそのまま提出すると、面接で自分の言葉で語れなくなる
\item 複数の就活生が同じAIを使うため、没個性的な文章になりやすい
\item AIは「あなたの経験」を知らない。事実の入力は必ず自分で行う必要がある
\item 企業によってはAI使用を検出する仕組みを導入している場合がある
\end{itemize}
\end{itembox}

本章で提唱するアプローチは、\textbf{「AIに書かせる」のではなく「AIと一緒に磨く」}というスタンスである。

\section{ガクチカとは何か}
\index{がくちか@ガクチカ}

\subsection*{「学生時代に力を入れたこと」の本質}

「ガクチカ」とは「学生時代に力を入れたこと」の略称で、ほぼすべての企業のESや面接で問われる定番質問である。ガクチカが問うているのは、\textbf{「何をやったか」ではなく「どのように取り組んだか」}である。

\subsection*{企業が見ているポイント}

\begin{table}[H]
\centering
\begin{tabular}{lll}
\toprule
ポイント & 説明 & 具体的に見ていること \\
\midrule
課題発見力 & 状況の中から問題を見つける力 & 何を課題と認識したか \\
行動力 & 自ら考えて動く力 & どのような工夫をしたか \\
成果 & 行動の結果として何が得られたか & 定量的・定性的な成果 \\
学び & 経験から何を学んだか & 自己成長、気づき \\
再現性 & 入社後も同様に活躍できるか & 学んだことをどう活かすか \\
\bottomrule
\end{tabular}
\end{table}

\subsection*{研究をガクチカにする場合の特有の課題}

\begin{enumerate}
\item \textbf{専門用語の壁}: 研究内容をそのまま書くと、専門外の採用担当者には理解できない
\item \textbf{「やらされた」感の払拭}: 与えられたテーマの中で、自分がどのような課題意識を持ち、どう工夫したかを明確にする必要がある
\item \textbf{成果の表現}: 研究の社会的意義や、その成果が何に繋がるのかまで含めて説明する必要がある
\item \textbf{個人の貢献の明確化}: 研究室での研究はチームで行うことも多く、自分の独自の貢献を明確にしなければならない
\end{enumerate}

\section{AIを使ったガクチカ作成ワークフロー}
\index{STAR法@STAR法}

ガクチカ作成を4つのステップに分け、各ステップでAIをどのように活用するかを解説する。

\subsection{Step 1: 素材の棚卸し --- AIに壁打ちしてもらう}

最初のステップは、自分の経験を漏れなく洗い出すことである。AIを壁打ち相手にすることで、忘れていた経験や見落としていた側面を引き出せる。

\begin{promptbox}
\# 役割\\
あなたは大学院生の就職活動を支援するキャリアカウンセラーです。

\# 依頼\\
私は修士課程の学生で、ガクチカの素材を探しています。以下の情報をもとに、私の経験を深掘りする質問を10個生成してください。

\# 私の情報
\begin{itemize}
\item 専攻: [例: 機械工学]
\item 研究テーマ: [例: 自動運転における安全性評価手法の研究]
\item 研究期間: [例: 2024年4月〜現在]
\item その他の活動: [例: TAの経験、学会発表1回、研究室のゼミ運営]
\end{itemize}

\# 質問の方向性: 研究の中での困難や工夫、チームワーク、自発的な取り組み、失敗と学び、スキルアップの努力に関する質問
\end{promptbox}

このステップでは、AIに文章を書かせるのではなく、AIに「質問してもらう」ことが重要である。

\subsection{Step 2: STAR法での構造化}

素材が見つかったら、STAR法\index{STAR法@STAR法}で構造化する。STAR法は、経験を説得力のある形で整理するためのフレームワークである。

\begin{table}[H]
\centering
\begin{tabular}{lll}
\toprule
要素 & 内容 & ガクチカにおける役割 \\
\midrule
\textbf{S}ituation（状況） & どんな環境・背景だったか & 読み手に場面をイメージさせる \\
\textbf{T}ask（課題） & 何が課題・目標だったか & 取り組みの意義を示す \\
\textbf{A}ction（行動） & 具体的に何をしたか & 自分の主体性・工夫を示す \\
\textbf{R}esult（結果） & どんな成果が得られたか & 行動の効果を証明する \\
\bottomrule
\end{tabular}
\end{table}

\subsection{Step 3: AIで文章化（400字程度）}

STAR構造ができたら、これを読みやすい文章に変換する。

\begin{promptbox}
\# 役割\\
あなたはESライティングの専門家です。

\# 依頼\\
以下のSTAR構造を、400字以内のガクチカ文章にまとめてください。

\# STAR構造\\
{[Step 2で作成したSTAR構造をここに貼る]}

\# 文章の条件
\begin{itemize}
\item 400字以内（厳守）
\item 最初の1文で「何に力を入れたか」が分かること
\item 専門用語は最小限にし、非専門家でも理解できること
\item 数字を使って具体性を出すこと
\item 最後に「学び」や「今後への展望」を含めること
\end{itemize}
\end{promptbox}

\begin{itembox}[l]{Before/Afterの例}
\textbf{Before（構造化前の素材メモ）}:

自動運転の安全性の研究をしています。GSNというのを使って安全性を評価する手法を提案しました。論文を50本くらい読みました。国際学会に通りました。

\textbf{After（AIと磨いた完成稿）}:

自動運転システムの安全性を論理的に証明する研究に力を入れました。修士研究として取り組んだのは、「なぜこのシステムは安全と言えるのか」を体系的に説明する枠組みの構築です。従来の安全性評価は、確認項目が断片的で全体像が見えにくいという課題がありました。私は50本以上の先行研究を精読して課題を整理した上で、安全性の論拠を木構造で可視化する手法と、システムの状態を継続的に記録する仕組みを組み合わせた新しいアプローチを考案しました。週1回の研究ミーティングでは、自ら論点を整理した資料を作成し、議論を主導しました。その結果、この成果をまとめた論文が安全性工学分野の国際学会に採択され、研究室初の海外口頭発表を実現しました。この経験を通じて、先行研究の徹底的な調査に基づく課題設定と、仮説を検証しながら新しい解決策を構築する力を身につけました。（350字）
\end{itembox}

\subsection{Step 4: AIで推敲・添削}

文章化ができたら、採用担当者の視点でAIに添削してもらう。

\begin{promptbox}
\# 役割\\
あなたはメーカーの採用担当者（技術系）です。年間500本以上のESを読んでいます。

\# 依頼\\
以下のガクチカ文章を、採用担当者の視点で添削してください。

\# 添削の観点
\begin{enumerate}
\item 【第一印象】最初の1文で興味を持てるか
\item 【課題の明確さ】何が課題だったかが明確か
\item 【行動の具体性】何をしたかが具体的に書かれているか
\item 【主体性】「自分が」やったことが伝わるか
\item 【成果】成果が具体的か（定量情報があるか）
\item 【再現性】入社後も活かせそうな能力が見えるか
\item 【読みやすさ】文章は読みやすいか、冗長な部分はないか
\end{enumerate}

\# 出力形式\\
各観点について、○/△/×の評価と、具体的な改善提案
\end{promptbox}

添削のサイクルは2〜3回繰り返すと効果的である。重要なのは、\textbf{AIの修正案をそのまま採用するのではなく、自分の言葉で書き直すこと}である。面接では必ずガクチカの深掘りがあり、自分の言葉で語れなければ見抜かれる。

\section{研究を非専門家に伝える技法}

\subsection*{エレベーターピッチ（30秒で研究を説明）}
\index{えれべーたーぴっち@エレベーターピッチ}

エレベーターピッチとは、エレベーターに乗っている30秒程度の時間で、自分の研究を相手に伝える技法である。

\begin{itembox}[l]{30秒エレベーターピッチの構造}
\begin{enumerate}
\item \textbf{一言で言うと}（1文）：研究を一言で表現
\item \textbf{なぜ重要か}（1〜2文）：社会的な文脈での意義
\item \textbf{何をしたか}（1〜2文）：アプローチの概要
\item \textbf{何が分かったか}（1文）：主要な成果
\end{enumerate}
\end{itembox}

\subsection*{技術的な内容の平易化テクニック}

\begin{table}[H]
\centering
\begin{tabular}{ll}
\toprule
専門的な表現 & 平易な表現（比喩） \\
\midrule
Goal Structuring Notation & 安全性の根拠を木構造で整理する図 \\
形式検証 & 数学的に「絶対に間違いがない」と証明すること \\
機械学習モデルの過学習 & 丸暗記したけど応用問題が解けない状態 \\
API & ソフトウェア同士が会話するための共通言語 \\
\bottomrule
\end{tabular}
\end{table}

\begin{itemize}
\item \textbf{テクニック1: 日常の比喩を使う} --- 専門概念を身近なものに例える
\item \textbf{テクニック2: 「何が嬉しいか」を先に言う} --- How（何をしたか）よりWhy/What（なぜ重要か、何が嬉しいか）を先に伝える
\item \textbf{テクニック3: 数字で具体化する} --- 「多くの先行研究」ではなく「50本以上の論文を精読」
\end{itemize}

\section{ES全体でのAI活用}
\index{えんとりーしーと@エントリーシート}

\subsection*{志望動機の作成}

志望動機は、ガクチカ以上に「自分の言葉」が求められる設問である。ポイントは、AIに「文章」ではなく「論理構造」を出力させることである。構造ができたら、自分の言葉で文章化する。

\subsection*{自己PRの作成}

自己PRはガクチカと役割が異なる。

\begin{itemize}
\item \textbf{ガクチカ}: 具体的なエピソードを通じて能力を示す（帰納的）
\item \textbf{自己PR}: 自分の強みを宣言し、エピソードで裏付ける（演繹的）
\end{itemize}

\subsection*{注意: AIに丸投げしない}

\textbf{ESは面接の「台本」である。}ESに書いた内容は、面接で必ず深掘りされる。AIを使って作成する場合でも、以下を心がけること。

\begin{enumerate}
\item \textbf{素材（事実・経験）は必ず自分で入力する} --- AIは知らないことを捏造する
\item \textbf{AIの出力を鵜呑みにしない} --- 自分の経験と合っているか確認する
\item \textbf{最終的な文章は自分の言葉に置き換える} --- 音読して違和感がないか確認する
\item \textbf{面接で話す練習をする} --- ESを見ずに、自分の言葉で3分間話せるか試す
\end{enumerate}

\section{面接準備でのAI活用}

\subsection*{想定質問の生成}

面接準備で最も効果的なAI活用は、想定質問の生成である。

\begin{promptbox}
\# 役割\\
あなたはメーカー（技術系）の面接官です。10年以上の面接経験があります。

\# 依頼\\
以下のESの内容を読み、面接で聞きそうな質問を15個生成してください。

\# 質問のカテゴリ
\begin{itemize}
\item ガクチカの深掘り（5問）
\item 志望動機の深掘り（3問）
\item 研究内容に関する質問（3問）
\item 価値観・人柄に関する質問（2問）
\item 逆質問のヒント（2問）
\end{itemize}

各質問に対して、質問文、この質問の意図、回答のポイントを記載してください。
\end{promptbox}

\subsection*{模擬面接（AIに面接官役をさせる）}

ChatGPTやClaudeなどの対話型AIを使って、模擬面接を行うことができる。1回につき1つの質問をし、回答に対してフィードバックを述べてから次の質問に進む形式が効果的である。

\subsection*{面接後の振り返り}

面接が終わった後の振り返りにもAIは活用できる。質問と自分の回答を記録し、各回答の論理構造、具体性、質問の意図への対応をAIに分析してもらう。

\vspace{5mm}
\begin{itembox}[l]{本章のキーポイント}
\begin{enumerate}
\item \textbf{就活文書は「AIに書かせる」のではなく「AIと一緒に磨く」} --- 素材は自分、構造化と推敲にAIを活用する
\item \textbf{STAR法は汎用的な構造化ツール} --- Situation→Task→Action→Resultの順で整理すれば説得力が増す
\item \textbf{非専門家への説明は「Why/What」から始める} --- 「何をしたか」より「なぜ重要か」を先に伝える
\item \textbf{ESは面接の台本} --- AIの出力をそのまま使わず、自分の言葉で語れるようにする
\item \textbf{AIを使った模擬面接は効果的だが、人間との練習も併用する}
\end{enumerate}
\end{itembox}

\section*{演習問題}

\begin{enumerate}
\item[\textbf{演習9-1}] \textbf{素材の棚卸し}\\
本章のStep 1のプロンプトを使い、AIに自分の経験を引き出す質問を生成してもらいなさい。生成された質問に対してすべて回答し、ガクチカの素材となりそうなエピソードを3つ選びなさい。

\item[\textbf{演習9-2}] \textbf{STAR法による構造化}\\
演習9-1で選んだエピソードの1つを、STAR法で構造化しなさい。AIを使ってもよいが、最終的な構造は自分の言葉で記述すること。Situation, Task, Action, Resultそれぞれについて、3〜5文で記述しなさい。

\item[\textbf{演習9-3}] \textbf{ガクチカ文章の作成と推敲}\\
演習9-2のSTAR構造をもとに、400字以内のガクチカ文章を作成しなさい。AIに下書きを作成させた後、Step 4の添削プロンプトで少なくとも2回の推敲サイクルを行い、Before（初稿）とAfter（最終稿）を比較して提出しなさい。

\item[\textbf{演習9-4}] \textbf{エレベーターピッチの作成}\\
自分の研究について、30秒で説明できるエレベーターピッチを作成しなさい。以下の2つのバージョンを作成すること。
\begin{itemize}
\item 版A: 同じ分野の研究者向け（専門用語使用可）
\item 版B: 文系の採用担当者向け（専門用語不使用）
\end{itemize}
実際に声に出して読み、30秒以内に収まるか確認しなさい。

\item[\textbf{演習9-5}] \textbf{模擬面接の実施}\\
本章のプロンプトを使い、AIと模擬面接を実施しなさい。8問以上のやり取りを行い、AIからのフィードバックをもとに、自分の回答の改善点を3つ挙げなさい。
\end{enumerate}
